% --------------------------------
% PROJECTS
% --------------------------------

\mediumtitle{Research collaborations}

\smallskip

\noindent Only collaborators outside King's College London (for projects started between 2014 and 2023) and the University of Turin (for projects started before 2014 and after 2023) are reported.


\bigskip

\smalltitle{Projects for which Alessia Visconti has scientific responsibility}

\begin{singletablelist}

	\stlist{2025 - 2026}{
		\textsc{Title:} Digital transition and artificial intelligence in the social and healthcare sector in the Province of Cuneo: dissemination, best practices, and possible developments \\
		& \textsc{Funder:} Cassa di Risparmio di Cuneo Foundation; \textsc{Budget:} \euro{58,000} \\
		& \textsc{Role:} co-PI \\
		& \textsc{Objectives:} Analysis of the diffusion and socio-economic impact of digital technologies in the social and healthcare sectors in the province of Cuneo, Italy. \\
		& \textsc{Scientific contribution:} Alessia Visconti co-led an initial review of the scientific and grey literature, co-designed two surveys, one aimed at citizens and the other at health and social care professionals, and will contribute to their quantitative and qualitative analysis. The project will result in the publication of a monograph as part of the \emph{``Quaderni della Fondazione CRC''} series. As part of this project, she co-supervises a PhD student (Emanuele Koumantakis) and a research assistant (Marco Vitturini). 
	}
	
	\stlist{2025 - 2026}{
		  \textsc{Title:} Read Assessment and Decision Support for ICU Readmission Prediction (READ-ICU) \\
		& \textsc{Funder:} Horizon Enfield -- 2nd Exchange Scheme Open call; \textsc{Budget:} \euro{43,200} (mobility allowance for three researchers) \\
		& \textsc{Collaborators:} National Technical University of Athens, Greece \\
		& \textsc{Role:} co-PI \\
		& \textsc{Objectives:} Developing an ensemble approach to predict the risk of ICU readmission. \\
		& \textsc{Scientific contribution:} Alessia Visconti co-leads the development of an approach based on the \emph{Bayesian Model Averaging Framework} to integrates statistical, classical machine learning, and deep learning models to improve the prediction of ICU readmission risk using publicly ICU datasets~\cite{Kou26}. The models were identified through a systematic review of the literature, which also included a meta-analysis~\cite{Kou25}. As part of this project, she co-supervises a PhD student (Emanuele Koumantakis). 
	}
	

	\stlist{2024 - 2026}{
		  \textsc{Title:} Exploring the potential role of gut microbiota in improving outcomes of chronic myeloid leukaemia patients \\
	   & \textsc{Funder:} ESH John Goldman Research Award; \textsc{Budget:} \euro{80,000}  \\
		& \textsc{Collaborators:} Italian Institute for Genomic Medicine, Italy \\
		& \textsc{Role:} Work package leader (data analysis) \\
		& \textsc{Objectives:} Analysis of the interactions between the gut microbiome and response to pharmacological treatments in patients with chronic myeloid leukaemia. \\
		& \textsc{Scientific contribution:} Alessia Visconti used unsupervised machine learning techniques and Bayesian statistics to identify gut microbiome profiles associated with a deep molecular response, supporting the hypothesis that the gut microbiome plays a role in modulating drug efficacy. 
	}
	
	\stlist{2022 - 2024}{
		  \textsc{Title:} A Collaborative Approach to the Borne Uterine Mapping Programme (BUMP) Feasibility Study \\
		& \textsc{Funder:} BORNE medical research charity; \textsc{Budget:} £500,000 \\
		& \textsc{Collaborators:} Imperial College London, UK, Newcastle University, UK, and University of Turin, Italy \\
		& \textsc{Role:} Work package leader(data analysis). \\
		& \textsc{Objectives:} Analysis of transcriptional profiles associated with the onset of spontaneous labour through the integration of scRNA-seq, snRNA-seq, and spatial transcriptomics data. \\
		& \textsc{Scientific contribution:} After conducting quality control and preliminary analyses, Alessia Visconti characterised the transcriptional profiles using statistical and \\
		&  computational approaches to \emph{(a)}~assess the reproducibility of gene expression networks in scRNA-seq and snRNA-seq experiments, \emph{(b)}~identify differentially abundant cell subpopulations in scRNA-seq and snRNA-seq experiments, and \emph{(c)}~identify differentially expressed genes in spontaneous compared with non-spontaneous labour. The study enabled a systematic assessment of the strengths and limitations of single-cell and single-nucleus technologies, showing how their integration with spatial transcriptomics improves the resolution and interpretation of the cellular and tissue mechanisms underlying labour. 
	} 
			
	\stlist{2021 - 2024}{ 
		  \textsc{Title:} Understanding phenotype and mechanisms of spontaneous preterm birth in sub-Saharan Africa (PRECISE-SPTB) \\
		& \textsc{Funder:} Medical Research Council (MRC); \textsc{Budget:} £458,80 \\
		& \textsc{Collaborators:} Aga Khan University, Kenya, MRC Unit The Gambia at LSHTM, The Gambia, and Eduardo Mondlane University, Mozambique \\
		& \textsc{Role:} Work Package Leader (Knowledge Transfer) \\
		& \textsc{Objectives:} Determining the epidemiological and contextual nature of spontaneous preterm birth in three sub–Saharan African countries (Kenya, Nairobi and Mozambique), while developing technical infrastructure and training research scientists. \\
		& \textsc{Scientific contribution:} Alessia Visconti designed and delivered workshops, rated as ``exceptional'' by participants, on basic and advanced skills (\emph{e.g.}, version control and collaboration using Git/GitHub, Bash programming, and metagenomic data analysis). 
		}
		
	\stlist{2020 - 2021}{ 
		  \textsc{Title:} A multi-omics study to dissect the role of the gut microbiome in IgA nephropathy risk \\
		& \textsc{Funder:} King's College London - Peking University Health Science Centre Joint Institute for Medical Research; \textsc{Budget:} £74,000 \\
		& \textsc{Collaborators:} Peking University, China \\
		& \textsc{Role:} Work package leader (data analysis). \\
		& \textsc{Objectives:} Identification of genes and microbes associated with IgAN or with IgA glycosylation. \\
		& \textsc{Scientific contribution:} Alessia Visconti conducted several association studies, including one of the first GWAS of IgA glycosylation~\cite{Vis24}.  %Responsible for the work package dealing with \emph{in silico} characterisation and validation, using bioinformatics models, of microbes associated with IgA nephropathy and/or IgA glycosylation profiles. \\
		}
			
		\stlist{2019}{
		\textsc{Title:} Studying the microbiome-disease connection \\
		& \textsc{Funder:} King's Undergraduate Research Fellowship (KURF); \textsc{Budget:} £2,000 (£1,000 for a student bursary and £1,000 for laboratory costs)\\
		& \textsc{Role:} PI  \\
			& \textsc{Objectives:} Exploratory analysis aimed at assessing the effect of chronic diseases on the taxonomic and functional composition of the gut microbiome.\\
			& \textsc{Scientific contribution:} Alessia Visconti supervised a summer internship (4 weeks, 150 hours), which aimed at generating preliminary results to explore a potential new research direction. %June–August 2019
		}
	
		\stlist{2018}{
		\textsc{Title:} Studying the microbiome-immune system interplay \\
		& \textsc{Funder:} King's Undergraduate Research Fellowship (KURF); \textsc{Budget:} £3,000 (£2,000 for a student bursary and £1,000 for laboratory costs)\\
		& \textsc{Role:} PI  \\
			& \textsc{Objectives:} Exploratory analysis aimed at assessing the interactions between the taxonomic and functional composition of the gut microbiome and immune-system composition (immunophenotypes). \\
			& \textsc{Scientific contribution:} Alessia Visconti supervised a summer internship (8 weeks, 300 hours), which aimed at generating preliminary results to explore a potential new research direction. %, June–August 2018
		}
	
	
		\stlist{2017}{
		\textsc{Title:} Metagenome-Wide Association Study of Lipid Profiles in Twins \\
		& \textsc{Funder:} King's Undergraduate Research Fellowship (KURF); \textsc{Budget:} £3,000 (£2,000 for a student bursary and £1,000 for laboratory costs) \\
		& \textsc{Role:} PI  \\
			& \textsc{Objectives:} Exploratory analysis aimed at assessing associations between the taxonomic and functional composition of the gut microbiome and systemic lipid profiles (LDL and HDL). \\
			& \textsc{Scientific contribution:} Alessia Visconti supervised a summer internship (8 weeks, 300 hours), which aimed at generating preliminary results to explore a potential new research direction. %, June–August 2017
		}
		
		
		
	\stlist{2016 - 2018}{ 
		  \textsc{Title:} A high-resolution map of copy number and structural variation in Qatari genomes and their contribution to quantitative traits and diseases \\
		& \textsc{Funder:} Qatar Foundation; \textsc{Budget:} £160,521  \\
		& \textsc{Collaborators:} Sidra Medical and Research Centre, Qatar \\
		& \textsc{Role:} Work package leader (data analysis). \\
		& \textsc{Objectives:} Development and validation of an \emph{ensemble} approach for the identification of structural variants (deletions, duplications, inversions, translocations, and insertions) and characterisation of the mechanisms through which these variants influence disease susceptibility. \\
		& \textsc{Scientific contribution:} Alessia Visconti was the scientific lead for the statistical and computational analyses in several publications~\cite{Ros21,Ros24,Ali26} and co-supervised a PhD student (Niccol\`o Rossi), continuing to do so throughout his postdoctoral period.
	}
		
\end{singletablelist}

% \bigskip

\smalltitle{Projects to which Alessia Visconti participates as researcher}

\medskip

\noindent  Alessia Visconti has a prominent role in all projects, as shown by the number of publications in which she is first/last author.

\smallskip

\begin{singletablelist}


	\stlist{2023 - Present}{
		\textsc{Title:} Repertor.IO: Incorporating Patient Preference Studies into Clinical Research and Decision Model \\
		& \textsc{Funder:} \emph{Nord Ovest Digitale e Sostenibile (L’ecosistema dell'innovazione}; funded by the Italian Ministry of University and Research (MUR) under the National Recovery and Resilience Plan (PNRR) \\
		& \textsc{Collaborators:} University of Padua, Italy\\
		& \textsc{Objectives:} Developing a platform for the collection, annotation, and meta-analysis of patients preference studies.\\
		& \textsc{Scientific contribution:} After exploring, through a survey, the use of patient preferences in clinical care and decision making~\cite{DiB24}, Alessia Visconti contributed to the development of the meta-analysis framework, to the activities leading to the creation of the start-up (see the \textsc{Technology Transfer and Public Engagement} section for details), and to the preparation of funding applications (\euro{28,000} from Finpiemonte). As part of this project, she co-supervises a PhD student (Matteo Bracco).
	}

	\stlist{2023 - 2026}{
		\textsc{Title:} Impact of plant diversity on cardiometabolic health \\
		& \textsc{Funder:} Chronic Disease Research Foundation (CDRF) \\
		& \textsc{Collaborators:} King's College London, UK \\
		& \textsc{Objectives:} Assess how habitual diet influences the gut microbiota and the metabolites it produces and, in turn, cardiometabolic disease risk.\\
		& \textsc{Research consumption:}  Alessia Visconti was responsible for the preparation of metagenomic data, and contributed to the statistical analysis of multi\emph{-omics} and dietary data in the study described in~\cite{Pop25}. 
	}

	% \stlist{2023-2027}{
	% 	\textsc{Title:} \emph{``PUZZLE: aPproccio integrato alla mUltimorbidit\`a: ricerca di base, traslaZionale, clinica e formaZione muLtidisciplinarE''}\\
	% 	% & \textsc{Funder:} \emph{Italian Minister of Education, University and Research (MIUR)}\\
	% 	& \textsc{Objectives:} Developing a novel model for the study, management, and care of patients with multi-morbidity.\\
	% 	& \textsc{Role:} Leading researcher in the application of Deep Learning approaches to identify multi-morbidity patterns from epidemiological data on several diseases (\emph{e.g.}, cardiometabolic diseases and cancer) and environmental exposures (\emph{e.g.}, smoking, alcohol consumption, physical activity). Within this project, she is supervising a master student (Triparna Day). \newline \newline \newline \\
	% 	&
	% }

	\stlist{2023 - 2025}{
		\textsc{Title:} TrustAlert \\
		& \textsc{Funder:} Fondazione Compagnia San Paolo e Fondazione Cassa Depositi e Prestiti \\
		& \textsc{Collaborators:} Bruno Kessler Foundation, Italy, Links Foundation, Italy, Innovo Group (now DedaNext), Italy, Local Health Authority Alba-Bra (ASL-CN2), Italy, Cottolengo Hospital, Italy\\
		& \textsc{Objectives:} Developing AI-based solutions for \emph{(a)}~near-real-time identification of health emergencies, \emph{(b)}~mapping vulnerable and multimorbid populations, and \emph{(c)}~crate a \emph{living lab} for micro and macro simulations.\\
		& \textsc{Research contributes:} Alessia Visconti was involved in \emph{(a)}~supporting the development of an event-detection system based on the GDELT platform for identifying and tracking the evolution of public health emergencies~\cite{DeL25}, \emph{(b)}~analysing administrative healthcare data to predict short-term readmission ($<90$ days) using deep learning approaches~\cite{Con24}, and \emph{(c)}~analysing administrative healthcare data to identify population multi-morbidity patterns, using both a Convolutional Autoencoder (ConvAE) architecture and a Top2Vec-based approach. As part of this project, she co-supervised two PhD students (Emanuele Koumantakis, Sara De Luca), four research assistants (Matteo Bracco, Giuseppe Martino, Emanuele Pietropaolo, Michela Bersia), and two MSc students (Francesca Rondinone, Lorenzo Miglietta). 
	}



	\stlist{2021 - 2022}{
		\textsc{Title:} Predicting Response to Immunotherapy for Melanoma with gut Microbiome and metabolomics - The PRIMM Study \\
	   & \textsc{Funder:}  Seerave Foundation \\
		& \textsc{Collaborators:} University Medical Center, Groningen, The Netherlands and University of Trento, Italy\\
		& \textsc{Objectives:} Use multi\emph{-omics} data to predict which patients with melanoma would benefit from the use of immune checkpoint inhibitors. \\
		& \textsc{Scientific contribution:} Alessia Visconti conducted the analysis for the identification of glyco-markers~\cite{Vis23}, and contributed to the analysis of proteomic data~\cite{Ros22}. As part of this project, Alessia Visconti co-supervised a PhD student (Karla Lee) and a postdoctoral researcher (Niccol\`o Rossi) on the statistical analysis of multi\emph{-omics} data. 
	}
	
	
	\stlist{2020 - 2022}{ 
		\textsc{Title:} ZOE COVID Study \\
		& \textsc{Funder:}  Department of Health and Social Care (DHSC) \\
		& \textsc{Collaborators:} ZOE Ltd (King's College London start-up),  Massachusetts General Hospital, US\\
		& \textsc{Objectives:} Monitor the SARS-CoV-2 pandemic in near real time to optimise resource allocation and assess the effectiveness of prevention and treatment measures. \\
		& \textsc{Scientific contribution:} Alessia Visconti collaborated within the ZOE COVID Study (now ZOE Health Study) to develop the pipeline for analysing data submitted daily by approximately 4 million users~\cite{Mur21}. She also led two studies on the cutaneous symptoms of the infection~\cite{Vis21, Vis22}, and collaborated as data analyst in several other studies~\cite{Men20,Lee20,Zaz20,Hop21,Wil21,Sud21}. 
	}


	\stlist{2019 - 2021}{
		\textsc{Title:} Dissecting the mechanisms of immune-mediated inflammation: a systems-immunology approach \\
		& \textsc{Funder:} Medical Research Council (MRC)\\
		% & \textsc{Collaborators:} \\
		& \textsc{Objectives:} Identifying targetable networks of circulating immune cells and cytokines involved in the normal immune-mediated inflammation process, and identify those disrupted in immune-mediated inflammatory diseases. \\
		& \textsc{Scientific contribution:} Alessia Visconti was the leading researcher for the bioinformatics analyses aiming at \emph{(a)}~the reverse engineering of immune cell co-expression networks and their involvement in a set of autoimmune diseases, and \emph{(b)}~the identification of genetic variations and microbes/metabolites responsible for the development of such diseases. As part of this project, she supervised a postdoctoral researcher (Niccol\`o Rossi).
	}
	
	
	\stlist{2019 - 2020}{ 
		\textsc{Title:} Identification of multi-disease drug targets through systems-immunology dissection of immune ageing \\
		& \textsc{Funder:} Sanofi iAwards Europe 2019 \\
		& \textsc{Objectives:}  Identify a set of immune biomarkers for the early identification of individuals at higher risk of accelerated immune ageing. \\
		& \textsc{Scientific contribution:} Alessia Visconti was the scientific lead for the statistical analysis, and supervised a postdoctoral researcher (Niccol\`o Rossi) on machine-learning topics. 
	}
	
	\stlist{2018 - 2020}{
		\textsc{Title:} Industry-sponsored research collaboration \\
		& \textsc{Funder:} Danone Nutricia Research \\
		& \textsc{Objectives:} Study the role of the gut microbiome in modulating the effects of yoghurt consumption on human health. \\
& \textsc{Scientific contribution:} Alessia Visconti collaborated to the statistical analysis, and was responsible for extracting sequence information from metagenomic data~\cite{LeR22}. 
	}
	
	
	% \stlist{2018-2024}{
	% 	\textsc{Title:} \emph{``TREatment of ATopic eczema (TREAT) Registry Taskforce''}\\
	% 	% & \textsc{Funder:} \emph{Chronic Disease Research Foundation}\\
	% 	& \textsc{Collaborators:} more than 30 centers over 13 countries (see the registry website for details \url{treat-registry-taskforce.org}) \\
	% 	& \textsc{Objectives:} studying the genetic and environmental basis of atopic dermatitis (consortium). \\
	% 	& \textsc{Role:} Leading bioinformatician for the TwinsUK cohort~\cite{Gro21,Bud22,Sta25}. \\
	% 	&
	% }

	\stlist{2016 - 2018}{
		\textsc{Title:} Gut microbiome modulation of fasting glucose homeostasis and postprandial glycaemic response in TwinsUK and PREDICT: towards personalised diet for healthy aging \\
		& \textsc{Funder:} Chronic Disease Research Foundation (CDRF)\\
		& \textsc{Collaborators:} University of Trento, Italy\\
		& \textsc{Objectives:} Studying the influence of the gut microbiome on cardiometabolic health. \\
		& \textsc{Scientific contribution:} Alessia Visconti \emph{(a)}~developed a tool for the reproducible analysis of metagenomic data ()\emph{YAMP})~\cite{Vis18b}, \emph{(b)} performed the bioinformatics analysis of metagenomics and metabolomics data~\cite{Vis19,Bar20}, \emph{(c)}~collaborated to several studies exploring the influence of the gut metagenome on human health~\cite{Lou23,Nog23a,Nog23b}. 
	}

	\stlist{2014 - 2016}{
		\textsc{Title:} An integrative genomics approach for non-invasive diagnostic biomarkers discovery in IgA nephropathy \\
		& \textsc{Funder:} Medical Research Council (MRC) \\
		& \textsc{Collaborators:} Imperial College London, UK\\
		& \textsc{Objectives:} Identify glycomics and genetic biomarkers for improving IgAN diagnosis and treatment.\\
		& \textsc{Scientific contribution:} Alessia Visconti applied statistical and bioinformatics approaches for studying the role of IgA and its glycosylation profiles in the development of IgA nephropathy~\cite{Lom16,Dot21,Vis24}.
	}
	
	
	\stlist{2013 - 2018}{
		\textsc{Title:} Genomic analysis of Type 2 Diabetes in Qatar, towards diabetes personalized medicine \\
		& \textsc{Funder:} Qatar Foundation\\
		& \textsc{Collaborators:} Weill Cornell Medicine-Qatar, Qatar\\
		& \textsc{Objectives:} Identifying the genetics and epigenetics basis of Type 2 Diabetes. \\
		& \textsc{Scientific contribution:} Alessia Visconti developed and implemented an approach for the association testing of quantitative traits in related individuals (\emph{PopPAnTe})~\cite{Vis16}, and conducted the bioinformatics analysis for~\cite{AlM15} and~\cite{Zag18}. 
	}
	
	

	\stlist{2013 - 2015}{
		\textsc{Title:} Senescence and melanoma -- An integrative systems biology approach to characterise the link between reduced biological senescence and melanoma susceptibility \\
		& \textsc{Funder:} British Skin Foundation \\
		& \textsc{Collaborators:} Hospital Clinic of Barcelona, Spain, QIMR Berghofer Medical Research Institute, Australia, Beijing Institute of Genomics, China, Institute of Cancer and Pathology, UK, Erasmus MC University Medical Center, The Netherlands, Women's Hospital, US\\
		& \textsc{Objectives:} Identifying the genetics basis of melanoma skin cancer and of its risk phenotypes. \\
		& \textsc{Scientific contribution:} Alessia Visconti applied statistical and bioinformatics approaches for studying melanoma, melanoma risk phenotypes, and their connection with ageing~\cite{Rib16,Pui16,Hys18,Vis18a,Duf17,Vis19a,Vis20,Lan20,San20}. As part of this project, she co-supervised a PhD student (Marianna Sanna).
	}


	\stlist{2012 - 2013}{
		\textsc{Title:} LIMPET -- Isotropic And Anisotropic Lipophilicity To Model Permeability Of New Therapeutic Peptides \\
		& \textsc{Funder:} Compagnia di San Paolo \\
		% & \textsc{Collaborators:} \\
		& \textsc{Objectives:} Applying machine Learning approaches to predicts peptides permeability on the basis of experimental non-conventional lipophilicity indexes. \\
		& \textsc{Scientific contribution:} Alessia Visconti \emph{(a)}~evaluated the ability of some combinations of descriptors/algorithms to find the best model to predict the lipophilicity of small peptides~\cite{Vis15a} , and \emph{(b)}~performed the bioinformatics analyses to assess the ability of more than 200 compounds to act as hydrogen-bond donors~\cite{Erm14}.
	}
	

	\stlist{2007 - 2011}{
		\textsc{Title:} BioBITs -- Developing white and green biotechnologies by converging platforms from biology and information technology towards metagenomics \\
		& \textsc{Funder:} Regione Piemonte, \emph{Bando Converging technologies 2007} \\
		& \textsc{Collaborators:} University of Eastern Piedmont, Italy, Consiglio Nazionale delle Ricerche - Istituto Protezione Piante, Centro Colture Sperimentali Valle d'Aosta Srl, Isagro Ricerca Srl, Etica Srl\\
		& \textsc{Objectives:} Identifying and characterising populations of uncultivable bacteria living inside a symbiotic, arbuscular mycorrhizal fungus.\\
		& \textsc{Scientific contribution:} Alessia Visconti contributed to the development of a modular framework for the analysis of metagenomics sequences, and was responsible for the co-clustering module~\cite{Bon11}. 
	}
	

	% \dtlist{2004-2009}{
	% 	\textsc{Title:} \emph{``Realizzazione di modelli informatici per la valorizzazione della qualit\`a e la tracciabilit\`a delle produzioni in specie da frutto coltivate in Piemonte''}\\
	% 	% & \textsc{Funder:} \emph{Regione Piemonte}\\
	% 	% & \textsc{Collaborators:} \\
	% 	& \textsc{Objectives:}  \\
	% 	& \textsc{Role:} Alessia Visconti contributed to the development of computational approaches for the classification and traceability of fruits produced in Piemonte.\\
	% 	&
	% }

\end{singletablelist}

% \newpage
%
% \smalltitle{Projects with industrial partners to which Alessia Visconti participates}
%
% \begin{singletablelist}
%
% 	\stlist{2020 - 2022}{
% 		\textsc{Partner:} ZOE Ltd \url{https://health-study.joinzoe.com/})\\
% 		& \textsc{Role:}  Alessia Visconti \emph{(a)}~developed the analysis pipeline for the daily data provided by more than four million users~\cite{Mur21}, \emph{(b)}~led two studies investigating skin manifestations of SARS-CoV-2~\cite{Vis21, Vis22}, and \emph{(c)}~performed the bioinformatics analysis for several other studies~\cite{Men20,Lee20,Zaz20,Hop21,Wil21,Sud21}.\\
% 		&
% 	}
%
% 	\stlist{2019 - 2020}{
% 		\textsc{Partner:} Sanofi (\url{https://www.sanofi.com}) (``Sanofi iAwards Europe 2019'') \\
% 		& \textsc{Role:} Alessia Visconti carried out part of the bioinformatics analyses, and supervised a postdoctoral researcher (Niccol\`o Rossi).\\
% 		&
% 	}
%
% % \end{singletablelist}
% %
% %
% % \begin{singletablelist}
%
% 	\stlist{2018 - 2020}{
% 	\textsc{Partner:} Danone Nutricia Research (\url{https://www.nutriciaresearch.com}) \\
% 	& \textsc{Role:} Alessia Visconti performed the metagenomic data analyses, as described in~\cite{LeR22}.
% 	}
%
% \end{singletablelist}


\smalltitle{Research activities not supported by direct funding}

\begin{singletablelist}
	
	\stlist{2025 - Present}{
		\textsc{Title:} Use of generative AI for teaching \\
		& \textsc{Role:} co-PI  \\
		& \textsc{Objectives:} Development of a virtual patient simulation environment through which medical students can practise in realistic clinical scenarios in a safe environment. \\
		& \textsc{Scientific contribution:} Alessia Visconti co-coordinates the design and development of an interactive chatbot based on a multi-agent architecture. As part of this project, she co-supervised two MSc students (Boris Turcan, Samuele Perrotta) and a medical student (Mohamed Ayadi).  %, which will shortly be the subject of a randomised clinical trial aimed at assessing its effectiveness in promoting the development of critical and context-appropriate communication skills among students. 
		}
	
	\stlist{2024 - Present}{
		\textsc{Title:} Multi\emph{-omics} data integration \\
		& \textsc{Role:} co-PI \\
		& \textsc{Objectives:} Development of new methodologies for patient stratification based on their multi\emph{-omics} profiles \\
		& \textsc{Scientific contribution:} Alessia Visconti co-coordinated the development of a novel three-step tool which successfully stratified both synthetic patients and those with chronic myeloid leukaemia~\cite{DeL26}. As part of this project, she co-supervises a PhD student (Sara De Luca). 
		}
		

\end{singletablelist}

\smalltitle{Participation in research groups}

\begin{singletablelist}
	
	% \dtlist{?? - Present}{\textbf{Member} of the working group for Health Behaviour in School-aged Children (HBSC) Italia, which studies health-related behaviours among school-aged children.
	% % (\url{https://www.epicentro.iss.it/hbsc/informazioni-generali}).
	% }

	\stlist{2024 - Present}{\textbf{Co-coordinator} of the Working Group on \emph{Machine learning in clinical research} of the Italian Society of Medical Statistics and Clinical Epidemiology (SISMEC). 
	}

	\stlist{2016 - 2022}{\textbf{Researcher} within the TREatment of ATopic eczema (TREAT) Registry Taskforce (a consortium involving more than 30 research institutes in 9 countries)\\
		& \textsc{Publications:}~\cite{Gro21,Bud22,Sta25}.
	}
	
	
\end{singletablelist}